\documentclass[11pt]{article}
\usepackage[a4paper,margin=25mm]{geometry}
\usepackage{amsmath,amssymb,amsthm,hyperref}
\newtheorem{theorem}{Theorem}
\newtheorem{lemma}[theorem]{Lemma}
\newcommand{\cost}{\operatorname{cost}}
\newcommand{\OPT}{\operatorname{OPT}}
\setlength{\emergencystretch}{3em}
\title{Nested preferred paths with splay auxiliaries}
\author{Research note, 2 October 2026}
\date{}
\begin{document}
\maketitle

\noindent\textbf{Scope and status.}
This note proves a splay-backend version of the local Cha-Cha
construction. It retains the secondary reference, its preferred-path
decomposition, its fixed-owner history, and the no-copy unfolding into
actual keys. It is a different algorithm from standard multi-splay,
which stores an entire primary preferred path in one auxiliary splay.
The argument is a hand proof. Its physical realization uses the
no-copy unfolding, rooted-hull transfer, and equal-frontier role-change
lemmas in the existing Cha-Cha upper manuscript. It has not been
Lean checked, externally reviewed, or implemented as a data structure.

\paragraph{Model.}
Use the same model as that manuscript: ordinary key-ordered BST pointer
moves and rotations are charged; every original query starts at the
physical root; nodes have $O(\log n)$ auxiliary bits; computation on
already read data is free; inaccessible saved pointers are not free
jumps. No uniform word-RAM running-time bound is claimed.

Let the primary reference have $n=2^H$ leaves, the actual ordered keys.
The secondary reference $B$ is the midpoint routing tree on depth labels
$0,\ldots,H$. Put
\[
 K=1+\operatorname{height}(B)=O(1+\log(H+2)),
 \qquad \lambda=1+\log_2 K.
\]
All initial primary and genuine secondary preferences point right.
Let $P$ be first-free primary alternations, $I_0\le n-1$ the primary
initialization discrepancy, and $F$ the exact initialized fixed-owner
secondary bit count. The outer emissions for a query are its wrong
primary depths in increasing order, followed by $H$. Consequently
there are $m+P+I_0$ emissions and $P+I_0$ nonfinal emissions.

\begin{theorem}[Splay auxiliaries in the nested construction]\label{thm:main}
In the stated physical realization and cost model, the algorithm below
serves every sequence $X$ at cost
\[
 O\bigl((m+P(X)+F(X)+n)\lambda\bigr).
\]
Thus the established fixed-owner count conversion and Cha-Cha
dynamic lower bound give
\[
 O\bigl((\OPT(X)+n)(1+\log\log\log(n+4))\bigr).
\]
This theorem concerns the nested algorithm defined here, not unchanged
standard multi-splay.
\end{theorem}

\section*{One persistent virtual forest}

A primary preferred path beginning at depth $a$ has an exit alphabet
$[a,H]$ and routing keys $a,\ldots,H-1$. Its full virtual routing tree
is ordered by these depth labels. Genuine secondary preferred paths
are represented by splay auxiliaries, with root marks separating them.
The restriction to a nonroot suffix $[a,H]$ also contains the structural
boundary spine $\beta_a$: the surviving secondary-reference intervals
with left endpoint below $a$. Keep this whole spine in one splay
auxiliary. It has at most $K$ nodes. All other surviving intervals are
genuine canonical subtrees with their existing fixed-owner histories.
No independent secondary preference bits are created for $\beta_a$.

Every ordinary auxiliary is a secondary preferred path and therefore
also has at most $K$ nodes. Splay rotations change its shape, not its
node set or its genuine preference bits.

For accounting, a virtual routing item at label $j$ on one primary
path can be identified by that path's actual primary separator at
depth $j$. This is an analytical identity; the no-copy unfolding may
assign that separator to the image of a different abstract routing
node. On a primary cut after depth $d$, the items with labels at most
$d$ stay with the prefix, the deeper items stay with the old suffix,
and the incoming suffix brings its own deeper items. Thus there are
exactly $n-1$ virtual routing items globally at every checkpoint, and
the same population persists through the intermediate cut and graft.

Define
\[
 \Phi=\sum_{v\ \text{virtual routing item}}\log_2 s(v),
\]
where $s(v)$ counts descendants reachable through unmarked auxiliary
edges. Structural boundary auxiliaries and genuine preferred-path
auxiliaries use the same formula. The potential does not depend on
which fixed owner supplies a secondary bit or on which physical node
is an item's current unfolding image. At all intermediate states used
below,
\[
 0\le\Phi\le(n-1)\log_2 K.                    \tag{1}
\]
There is one global potential; no new initial potential is assigned
when a suffix changes status or is reattached.

\section*{Operations on the short auxiliaries}

The ordinary access lemma bounds a splay inside an auxiliary of size
at most $K$ by $O(1+\log K)$ amortized cost. This includes a splay
stopped at a supplied auxiliary subtree root. Searches are paid by
splaying their last auxiliary node; their pointer walk, return, and
splay cost differ by only a constant factor. Use a fixed constant
multiple of $\Phi$ to account for the chosen pointer convention.

The usual preferred-child switch is available in this same virtual
BST. Splay its reference vertex $y$ to its auxiliary root. Splay its
nearest reference-ancestor predecessor and successor, when present
in that auxiliary, below $y$. This exposes the old and incoming
tails in the corresponding gaps. Mark the old tail before unmarking
the incoming tail. These are the familiar three-splay switch operations.
Marking cannot increase $\Phi$. Unmarking increases auxiliary sizes
only at $y$ and at most one exposed boundary child, hence increases
$\Phi$ by at most $2\log K$. The resulting component is precisely
the new secondary preferred path, whose size is at most $K$.

The switch bound is local and does not require the bottom-up scheduling
used in standard multi-splay. Here use it top-down: search the current
root auxiliary toward the requested depth leaf $d$. If the search
enters a marked auxiliary or exits to a nonpreferred external reference
leaf, splay the last node of the old root auxiliary to pay for that
walk, then switch the first wrong secondary preference and repeat
from the virtual root. The latter case matters when the incoming
preferred path has no routing nodes and hence no marked root. The
stored bottom leaf of the root reference path and the fixed $B$
distinguish it from successful completion without a pointer jump.
The switched reference
vertex and its boundary keys are determined by the fixed reference
and the encountered exit. Finding them is by ordinary paid searches
and splays. The already read topology specifies whether a boundary
key is present, so an unsuccessful full scan is not hidden.

When the currently preferred depth leaf $d$ is reached, splay the last
routing node to pay for the final walk. The whole root-to-$d$ reference path is now one
auxiliary. If the inner request changes $f$ genuine bits, this procedure
costs $O((f+1)\lambda)$ amortized. It writes exactly the genuine bits
on the reference path to $d$ and leaves all other histories intact.
Only the current root primary path is queried in this manner, so the
secondary reference during each request is the full $B$.

\section*{The primary splice and its exact potential account}

Consider a nonfinal outer emission $d<H$ after the preceding inner
request to depth leaf $d$. Write $A$ for the root auxiliary. It has
at most $K$ items and contains the whole $B$-path to leaf $d$.
In any full routing tree, the two boundaries adjacent to an interior
leaf are its ancestors. Thus routing key $d$ belongs to $A$, and so
does $d-1$ when $d>0$.

\paragraph{1. Expose and detach.}
Splay routing key $d$ to the root of $A$. Split off its right virtual
subtree, retaining $d$ with the prefix. The root auxiliary's keys
greater than $d$ are exactly $\beta_{d+1}$, by the suffix invariant.
Their genuine hanging subtrees are unchanged. They become the old
suffix's existing structural boundary auxiliary, rather than a newly
initialized backend. If $\beta_{d+1}$ is empty, the suffix starts at
an already marked genuine canonical subtree, or is a single leaf.

The split only removes auxiliary descendants from the prefix root,
so its potential change is nonpositive. The suffix items keep their
internal auxiliary sizes. Splay $d-1$, when it exists, inside the left
subtree of $d$. This costs $O(\lambda)$ amortized. Now the last two
prefix depth leaves, $d$ and its terminal hole $d+1$, are incident
to the rooted abstract hull consisting of $d$ and $d-1$ (only $d$
when $d=0$).
Visit the actual frontier root at exit $d$ and read the incoming
primary-path header there. The exposed hull makes this an $O(1)$
physical walk. If that block is one actual key, visit the key and
use the empty suffix state instead. This is an actual pointer visit,
not an access to a saved but inaccessible header.

\paragraph{2. Exchange exit roles.}
Exchange the roles of the old depth-$d$ off-path block and the terminal
hole. These are adjacent actual-key intervals separated by the actual
primary separator at depth $d$. Their physical key order is unchanged.
Every other frontier component is unchanged. The exposed abstract hull
has at most two nodes, so its unfolding has constant size; the
equal-frontier role-change lemma realizes the exchange in $O(1)$
physical pointer operations and rotations. The virtual prefix topology,
auxiliary marks, and auxiliary sizes are unchanged. Consequently
\[
 \Delta\Phi_{\mathrm{role\ exchange}}=0.            \tag{2}
\]
This is why exposing $d-1$ was necessary: merely knowing the two exit
labels would not supply their physical interface in an arbitrary
splay shape.

\paragraph{3. Graft and concatenate.}
Graft the incoming primary suffix into the new terminal hole. Its
stored $\beta_{d+1}$ and all genuine hanging subtrees are preserved.
Since $d$ is the maximum routing key of the prefix and is its root,
the incoming virtual routing tree is attached at its right gap.
If the incoming boundary auxiliary is nonempty, clear that boundary
root's auxiliary mark. This concatenates its keys, all greater than
$d$, with the prefix auxiliary. No traversal down the incoming
boundary spine is required and no additional splay is necessary.
Only $d$'s auxiliary size increases. If the old prefix size is $p$
and the incoming boundary size is $b$, then
\[
 \Delta\Phi_{\mathrm{concatenate}}
 =\log_2\frac{p+b}{p}\le\log_2 K.                \tag{3}
\]
Here $p+b\le K$: the resulting auxiliary is exactly the full reference
path toward leaf $d$. If the incoming boundary is empty, the genuine
root mark remains and (3) is zero.

Outer primary root marks are distinct from these inner auxiliary marks.
Clearing the incoming outer mark and marking the parked old tail does
not change $\Phi$. Their physical roots are frontier roots of the
supplied exposed hull and can be visited at constant additional cost.
A leaf tail requires no root or topology header.

\paragraph{Which histories survived?}
Every genuine interval whose owner depth exceeds $d$ travels with its
incoming or parked suffix. An interval whose owner depth is at most
$d$ retains the prefix owner's bit. Any interval crossing the splice
contains $d$, and its bit was just set toward $d$ by the inner request.
In particular, the incoming boundary nodes have split greater than
$d$ and become genuine members of the root auxiliary with their
prefix-owned bit set left, exactly as required. The physical boundary
items did not carry an independently reset history. Reclassifying them
as genuine does not change their ranks in $\Phi$; their only potential
change is the explicitly charged unmarking in (3). The old suffix's
deeper-owner bits are preserved and none is reset.

Thus the splice costs $O(\lambda)$ amortized with respect to the same
global potential. There is no size-of-suffix startup payment.

\section*{Physical simulation, accounting, and the distinction from multi-splay}

The virtual operations above use actual depth-key comparisons. Their
physical realization is the no-copy unfolding: each abstract interval
has at most two actual-key components, each separator occurs physically
once, and a rooted touched hull of $q$ abstract nodes can be reconfigured
in $O(q)$ physical BST cost while preserving its frontier components.
Splay rotations, searches, and the supplied split hulls meet that
contract. A whole root-starting virtual transaction is rendered from
its touched hull, so one does not restart a physical root search for
each abstract rotation. The role exchange separately uses the constant
hull already exposed above.

Read the current root primary-path header, which includes its bottom
leaf $z$. Before the inner request, compute the outer exit from these
already read data: $d=H$ if $x=z$, and otherwise
$d=\operatorname{depth}(\operatorname{LCA}_T(x,z))$. This is allowed
computation on the query and the fixed reference; it does not read an
unvisited incoming header. The inner request and, for $d<H$, the two
boundary splays above then expose the actual exit. Its constant-length
physical walk reads the incoming header as specified in step 1.

For the final exit $d=H$, finish the inner request, then splay routing
key $H-1$ to its auxiliary root, at $O(\lambda)$ amortized cost.
It belongs to the exposed full reference path. Its right abstract
frontier is leaf $H$, so the constant-size unfolding hull gives an
$O(1)$ physical walk to the requested actual key $x=z$; perform this
walk before finishing. For an empty routing universe, visit the sole
actual key directly. Thus the procedure really visits every requested
key and does not end after only a header computation.

Incoming and parked primary headers are read or
installed only at actually reached physical frontier roots. All
virtual topology, marks, and owner bits of a primary path require
$O(H)$ bits, so fit in a constant number of $O(\log n)$-bit words.
This supplies the bounded augmentation in the stated free-computation
model, not a claim that decoding runs in constant word-RAM time.

The outer procedure fixes the first wrong primary preference before
continuing deeper. Hence its emissions are exactly the prescribed
increasing-depth stream. During an inner access to $d$, every visited
secondary interval starts at a depth $a\le d$, where the current root
primary path agrees with the requested key. Its owner is therefore
the fixed primary ancestor at depth $a$. The total number of genuine
secondary changes is exactly $F$, independent of splay shapes.

Summing the inner-request and splice bounds, with one fixed constant
multiple of $\Phi$, gives
\[
 \cost(X)+c\bigl(\Phi_{\mathrm{final}}-\Phi_{\mathrm{initial}}\bigr)
 \le C\lambda\bigl(m+P+I_0+F\bigr).
\]
By (1), the potential telescope contributes at most $O(n\lambda)$.
The initial legal physical shape and its headers can be installed in
$O(n)$ cost from any prescribed starting BST. Using $I_0\le n-1$
proves the first bound in Theorem~\ref{thm:main}. The separate count
inequalities $F=O(D+P+n)$ and $m+P+D=O(\OPT)$ give its competitive
corollary.

The extra secondary decomposition is essential to this proof: it caps
every splay auxiliary at $K=O(\log H)$ instead of $H$. Standard
multi-splay supplies only the latter cap. Neither the per-switch
access-lemma calculation nor the splice telescope just proved bounds
standard multi-splay by this nested algorithm's cost.

\paragraph{Adopted sources.}
The physical unfolding, suffix invariant, and exact-owner construction
are Sections 3--6 of the sibling Cha-Cha upper-bound manuscript,
\path{research/chacha-upper-bound.tex}.
The three-splay preferred-child switch and its sum-log-size accounting
are described in Wang's thesis, Sections 3.2--3.4.1,
\url{https://www.csd.cmu.edu/sites/default/files/phd-thesis/CMU-CS-06-140.pdf}.
The proof here supplies the additional amortized boundary exposure,
exact unchanged-potential role exchange, and single global splice
potential needed for the splay-backend substitution.

\end{document}
